% =====================================================================
% Response to the supervisor's \federico{} comments on Paper I and Paper D.
% Standalone: compiles with a bare `pdflatex response_supervisor_I_D.tex`.
% =====================================================================
\documentclass[11pt,a4paper]{article}

\usepackage[T1]{fontenc}
\usepackage[utf8]{inputenc}
\usepackage[expansion=false]{microtype}
\usepackage[margin=2.4cm]{geometry}
\usepackage{amsmath,amssymb}
\usepackage{booktabs}
\usepackage{enumitem}
\usepackage{xcolor}
\usepackage{titlesec}
\usepackage[hidelinks]{hyperref}

\definecolor{qgrey}{HTML}{4A4A4A}
\definecolor{rule}{HTML}{C9C9C4}
\definecolor{tag}{HTML}{8A4B2A}

\titleformat{\section}{\large\bfseries}{\thesection}{0.7em}{}
\titlespacing*{\section}{0pt}{1.6em}{0.7em}
\setlength{\parskip}{0.35em}
\setlength{\parindent}{0pt}

\newcommand{\thinrule}{\textcolor{rule}{\rule{\linewidth}{0.4pt}}}
\newenvironment{qitem}[3]{%
  \par\medskip\thinrule\par\smallskip
  \noindent\textbf{\textcolor{tag}{#1}}\quad\textcolor{qgrey}{\small #2}\par\smallskip
  \noindent\textcolor{qgrey}{\small\itshape #3}\par\smallskip
}{\par}
\newcommand{\where}[1]{\par\smallskip{\small\textcolor{tag}{$\hookrightarrow$}\ \textit{#1}}\par}

\newcommand{\Dauth}{\Delta_{\mathrm{auth}}}
\newcommand{\gu}{\gamma_{u}}
\newcommand{\ga}{\gamma_{a}}
\newcommand{\Rtx}{R_{\mathrm{tx}}}
\newcommand{\tfresh}{\tau_{\mathrm{fresh}}}
\newcommand{\eclk}{e_{\mathrm{clk}}}
\newcommand{\Norm}[1]{\lVert #1\rVert}

\begin{document}

\begin{center}
{\Large\bfseries Response to Supervisor Comments}\\[0.35em]
{\large Paper I (authentication budget) and Paper D (cross-layer certificate)}\\[0.9em]
\begin{tabular}{@{}ll@{}}
\textbf{Re} & 15 comments on Paper I and 7 comments on Paper D\\
\textbf{Outcome} & All addressed; both papers then received a full language pass\\
\end{tabular}
\end{center}

\section*{Summary}

Thank you for the careful reading. Each comment is answered below, with the place
in the revised manuscript where the change was made. Three points summarise the
revision.

\textbf{Paper D now uses one value of $\gamma$ throughout.} The supremum rate
$\gamma=1.625$\,m/s is the only rate used after \S6.3. Section~9 and Figures~6--7
were regenerated from the result files at that rate. The arithmetic you flagged in
\S6.3.3 was wrong: the tube at $\theta=0.25$\,s is $2.65$\,m, not $5.9$\,m.

\textbf{Paper I no longer depends on Paper D.} All references to ``the parent'' and
all citations of the unpublished manuscript have been removed. The certificate that
Paper I uses is now stated and proved inside Paper I (\S II-A and the appendix).

\textbf{The experiment list in Paper I now matches what was run.} The numbered
experiments E1--E8 are gone. A new table, \emph{Validation Status of the Inputs}, states, input by input, which checks were performed and which are
specified but not yet run.

\textbf{Language.} Following your general comment, both papers have been reread
and edited for tone and register: rhetorical devices, emphatic adverbs, aphorisms
and question-style openings were removed, and spelling was made consistently
British. Mathematics, numbers, labels and citations were checked mechanically to be
unchanged by this pass.

% ---------------------------------------------------------------------
\section{Paper I}

\begin{qitem}{I-1}{Abstract, finding (ii)}{The 1\,cm figure is the difference between
the two error envelopes before $\Dauth$ is paid; the net gain at $\theta=0.25$\,s is
slightly negative. Calling 1\,cm the gain contradicts \S V-C and Table V.}
Agreed. Finding (ii) now reads: ``At detector-admitted delays the two error
envelopes differ by 1\,cm, within noise, and the net effect is slightly negative;
authentication pays only from $\theta_\times\approx0.259$--$0.284$\,s.'' The
abstract remains at 249 words.
\where{Abstract.}
\end{qitem}

\begin{qitem}{I-2}{Introduction and throughout}{Stop calling the other paper ``the
parent''; it is unpublished and cannot be cited or posted without breaking
anonymity.}
Done. Every occurrence of ``parent'' and every \verb|\cite{parentpaper}| has been
removed, including the row in the related-work table and the node in Fig.~1. The
delay certificate is now introduced in a new background subsection, ``A
Certificate for Delayed Navigation Input'' (\S II-A). It states the residual-delay
lemma and the tube lemma, and both are proved in the appendix, so Paper I is
self-contained. The overlap with Paper D is limited to that background; if both
papers go to the same venue, we will declare it to the programme chairs as the
concurrent-submission policy requires.
\where{\S I, \S II-A (new), Appendix (two new proofs), Table of related work, Fig.~1.}
\end{qitem}

\begin{qitem}{I-3}{Introduction}{Did you run all the experiments? You mention eight.}
No, and the paper now says so plainly. Two checks were performed: the
response-time bound against a discrete-event simulation, and the GNSS-outage
dead-reckoning experiment (together with the decoding of the HCRL capture). The
others (SITL spoofer sweep, on-target cycle counts, raw-bus timing, GNSS time
push) are specified but not yet run. The E-numbering has been removed. A new
subsection and table, ``Validation Status of the Inputs'', mark each input as
performed or not run, and Limitations refers to it.
\where{\S IV (new subsection and table), Limitations, Fig.~1 box (5).}
\end{qitem}

\begin{qitem}{I-4}{All headings}{Revise all section, subsection and paragraph titles.}
Done. Headings are now descriptive noun phrases (for example ``Crossover under a
Constant Rate'', ``Transmission Cost and the Admissible Rate'', ``RQ2: Cost of
Authentication'', ``Certification versus design'').
\where{Throughout.}
\end{qitem}

\begin{qitem}{I-5}{\S II}{Clarify the physical path linking $\Delta(\theta)$ (mesh),
$\Dauth$ (onboard bus) and $\gu\to\ga$ (RF flights).}
A paragraph, ``Physical path of a navigation message'', now follows a correction
from the peer mesh, where it may be delayed by up to $\Delta(\theta)$, to the
onboard bus, where it is authenticated at cost $\Dauth$, and into the estimator.
The two delays act in series and add. The rates $\gu$ and $\ga$ are measured at
the estimator input, which is where both delays take effect. The transfer of
those rates from RF-attack flights to delayed mesh traffic is stated as an
explicit assumption (Assumption~1, ``Interface transfer'').
\where{\S II-A.}
\end{qitem}

\begin{qitem}{I-6}{Threat model}{Table~III lists six adversaries, not two.}
Corrected. The text now says there are six adversaries. The delay and denial
adversaries define the two regimes, and the other four arise only once
authentication is added.
\where{\S II, threat model.}
\end{qitem}

\begin{qitem}{I-7}{Proposition on GNSS time push}{``Experiment E7 tests it'' cannot sit
inside a proposition.}
Removed from the proposition. A remark after the proposition refers to the
validation table, where the GNSS time-push experiment is listed as not yet run.
\where{\S III, Proposition and following remark.}
\end{qitem}

\begin{qitem}{I-8}{Eq.~(amort-lat)}{Write $\Rtx(k)$, since the transmission response
time depends on the batch size.}
Done, and $\Rtx(k)$ is now defined in the text. The text also states that $n(0)=0$,
so no path frames are sent when $k=1$.
\where{\S III, Batching and Configuration Selection.}
\end{qitem}

\begin{qitem}{I-9}{Algorithm~1}{The regime check in Line~2 is never referenced again.}
A sentence before the algorithm now explains how the regime is used. Under denial,
unauthenticated sources are withheld beyond $T_c$, so the baseline (none, 0) is
inadmissible; Line~2 removes it, and Lines~5--16 then compare only the
authenticated configurations.
\where{\S III, text before Algorithm~1.}
\end{qitem}

\begin{qitem}{I-10}{Provenance tags}{``Four tags are used'' followed by five.}
Corrected to five, with a note on the \textsc{proxy} qualifier.
\where{\S IV, Inputs and Their Provenance.}
\end{qitem}

\begin{qitem}{I-11}{Transport}{Write out the frame-count formula.}
Added: $n(\sigma)=\lceil(\sigma+2)/7\rceil$ on classic DroneCAN (seven payload
bytes per frame plus a two-byte transfer CRC) and $\lceil(\sigma+2)/63\rceil$ on
CAN~FD, with $n(0)=0$. We checked the formula against the frame column of the
scheme table (MAC~3, Ed25519~10, FN-DSA~96, ML-DSA~346, SLH-DSA-s~1123,
SLH-DSA-f~2442, TESLA~4) and against the artifact's transport model.
\where{\S IV, Transport paragraph.}
\end{qitem}

\begin{qitem}{I-12}{Configuration grid}{Write ``two post-quantum flag settings''.}
Done.
\where{\S IV, Configuration Grid.}
\end{qitem}

\begin{qitem}{I-13}{RQ1}{The headline ratio uses \texttt{\textbackslash nRatioHead}
in the abstract and \texttt{\textbackslash nGammaRatio} here.}
The two macros hold the same value (1.74). RQ1 now uses \verb|\nRatioHead|, so the
abstract, \S III and RQ1 share one macro.
\where{\S V, RQ1.}
\end{qitem}

\begin{qitem}{I-14}{Freshness table}{Clarify what the 0.25\,s margin is compared against.}
The text and the table note now say explicitly that the margin is the slack
$\Delta(0.25)-\tfresh=0.50-0.25=0.25$\,s left for clock error $\eclk(t)$ before the
binding condition $\tfresh+\eclk(t)<\Delta(\theta)$ fails.
\where{\S V, RQ4 and Table (freshness).}
\end{qitem}

\begin{qitem}{I-15}{Figure (flooding)}{This figure is not referenced in the text.}
It is now referenced in RQ5.
\where{\S V, RQ5.}
\end{qitem}

% ---------------------------------------------------------------------
\section{Paper D}

\begin{qitem}{D-1}{Notation}{$m$ and $\delta$ are used for different things; $\theta$
and $\varepsilon$ in the appendix; $LT$ and $LT_c$.}
Fixed. $m$ now denotes only the corridor margin; the PAC-Bayes sample size is
$n_{\mathcal{M}}$ and its confidence parameter is $\beta$. $\delta u$ denotes the
input perturbation and $\delta(\theta)$ only its bound; the two places where
``$\delta$'' stood for the calibration now say $\gamma$. The detector threshold is
$\theta$ everywhere, including the operating-curve axes and Appendix~B, and the
notation table no longer lists $\varepsilon$ as an alias. Every exponent is
$LT_c$. The MWU regret in Table~V now uses $N$ for the number of rounds, as the
text does.
\where{\S6.3, \S8, \S9, Appendices A--B, Tables V and notation.}
\end{qitem}

\begin{qitem}{D-2}{\S6.2 / \S6.3.2}{State the physical link between forwarding delay and
a $\gamma$ measured under jamming and spoofing.}
Added to \S6.2. A delayed peer correction leaves the estimator without fresh
information, which is physically the jamming (coasting) regime, where $\gamma$ is
lowest. We nevertheless certify at the spoofing supremum. This is conservative for
pure staleness, and it remains a bound if a delayed report passes the innovation
gate and pulls the estimate as a counterfeit fix would. A sentence in \S6.3.2
points back to this argument.
\where{\S6.2 (new paragraph), \S6.3.2.}
\end{qitem}

\begin{qitem}{D-3}{\S6.3.3, \S6.3.4, \S9.2--9.5, Figs.~6--7}{One $\gamma$ throughout, and
the \S6.3.3 arithmetic.}
You were right on both counts. With $\gamma=1.625$\,m/s, $H=2$, $\theta=0.25$\,s and
$L=1.181$, $\rho=0.50\times1.625\times e^{1.181}=2.65$\,m. \S6.3.3 now says so,
and it states that $1.625$\,m/s is the only rate used in the rest of the paper.
The secant values ($1.12$\,s, $2.22$\,m) now appear only once, as the comparison
that motivates the choice. Every dependent number was recomputed from the result
files:
\begin{itemize}[leftmargin=1.4em,itemsep=1pt]
\item window edge at $m=10$\,m: $\theta^*=0.94$\,s (window $[0.178,0.944]$\,s);
\item worked example: $\Norm{\delta u}\le0.81$\,m, $\rho\approx2.65$\,m, floor
$0.75$ (only the 2\,m corridor fails); kinematic tube $24.4$\,m;
\item $\gamma_{\mathrm{req}}=6.14$\,m/s is $3.8\times$ the measured rate (was
``4.5--5.1$\times$'');
\item heading error needed for cruise to cross $\gamma_{\mathrm{req}}$: $0.30$\,rad,
about $17^\circ$;
\item the linear-model domain argument in \S6.3.4 now compares
$\tau_{\mathrm{lin}}\approx2$\,s with $0.94$\,s.
\end{itemize}
Figures~6 and~7 and the pairing and constants tables were regenerated by the
figure script, which now reads the supremum from the per-sample campaign file. The
$2.22$ versus $2.25$\,m discrepancy is gone, because both now read $2.65$\,m. A
further stale value turned up during this check: the kinematic window was quoted
as $[0.178,0.243]$\,s, which came from the global Lipschitz constant. At the local
constant used everywhere else, that window is empty at 10\,m and is
$[0.178,0.205]$\,s at 20\,m. The text now says this.
\where{\S6.3.3, \S6.3.4, \S9.2--9.5, \S10, Figs.~6--7, Tables V and X.}
\end{qitem}

\begin{qitem}{D-4}{\S5.1, Table~IV, \S9.6, \S11}{Contradictions about
\texttt{measured\_same\_platform} pairings.}
Reconciled along the lines you suggested; we did not annotate attack windows by
hand. All four passages, and the corpus table caption, now say the same thing: the
UAV Attack Dataset supplies same-platform \emph{calibration flights} for $\gamma$,
and because its release lacks machine-readable attack intervals, the artifact
contains no \texttt{measured\_same\_platform} pairing records (\S9.7). The pairing
table's \texttt{pairing\_basis} row now lists the two classes that actually occur.
\where{\S5 (Table I caption, \S5.1), Table IV caption, \S9.6, \S11.}
\end{qitem}

\begin{qitem}{D-5}{Appendix D, Definition 2}{The Grönwall step does not give
$\delta(e^{LT_c}-1)$; Definition~2 and $e(0)=0$ disagree.}
Agreed; the step was wrong. The proof now applies the comparison lemma to
$\tfrac{d^+}{dt}\Norm{e}\le L\Norm{e}+L\delta(\theta)$ with $e(0)=0$, giving
$\delta(\theta)(e^{LT_c}-1)$, and shows that the full integral form of Grönwall's
inequality yields the same result. Definition~2 now states that the margin $m$ is
defined net of the filter's nominal estimation error, so the deviation starts at
zero at each re-anchoring. The proof also records your second reading: if
re-anchoring leaves an offset $\Norm{e(0)}\le\delta(\theta)$, Grönwall gives
$\delta(\theta)e^{LT_c}$ directly. The reported radius $\delta(\theta)e^{LT_c}$ is
therefore valid under both readings, and no number changes. The theorem now assumes
$F$ is Lipschitz in state \emph{and} input, which the proof uses.
\where{Definition~2, Theorem~1 and its proof sketch, Appendix~D.}
\end{qitem}

\begin{qitem}{D-6}{Eq.~(floor), Fig.~6}{Clarify how the empirical probability is
computed.}
A paragraph, ``How the floor is computed'', now opens \S9.4, and the Fig.~6
caption repeats it in brief. For each $\theta$, $H$ is fixed by the topology
($H=2$ on the grid), so $\rho(\theta)$ is a single number. The floor is the
equally weighted fraction of the margin family $\{2,5,10,20\}$\,m with
$\rho(\theta)\le m$, so it takes values in $\{0,0.25,0.5,0.75,1\}$. The temporal
conjunct is not evaluated numerically, because the navigation anchor records no
mission durations. The paragraph says this and gives the mission length above which
the conjunct holds ($5$\,s at $\theta=0.25$\,s; $100$\,s at saturation).
\where{\S9.4, Fig.~6 caption.}
\end{qitem}

\begin{qitem}{D-7}{Appendix A}{``A reader fluent in \dots\ can proceed directly'' makes no
sense in an appendix.}
Removed.
\where{Appendix A.}
\end{qitem}

% ---------------------------------------------------------------------
\section{Further corrections found during the revision}

Rereading both papers line by line turned up errors that you had not marked. We
corrected them and list them here so that none of the changes is silent.

\textbf{Paper D.}
\begin{itemize}[leftmargin=1.4em,itemsep=2pt]
\item \emph{Linear domain of $\gamma$ (\S6.3.4).} The text said that the whole
certified window lies inside the linear model's domain ($\tau_{\mathrm{lin}}\approx2$\,s).
That compared $\theta^*$ with $\tau_{\mathrm{lin}}$, whereas the condition applies to
$\Delta(\theta)=H\theta$. v8 was already wrong here: $2\times1.12=2.24$\,s exceeds
$2$\,s. The text now says that the operating point ($\Delta=0.50$\,s) and the
$10$\,m window edge ($\Delta=1.89$\,s) lie inside the domain, but that for
$\theta>1$\,s, including the $20$\,m window edge, the linear model is applied
outside its domain. Values in that range are therefore upper bounds, not certified
values, because $b_{\mathrm{IMU}}$ has not been measured. No result at the
operating point is affected.
\item \emph{Significance (\S9.5).} ``Every adjusted $p<10^{-20}$'' held only for an
earlier mapping. At $\gamma=1.625$\,m/s the largest Holm-adjusted $p$ is
$7.9\times10^{-14}$, so the text now says $p<10^{-13}$. It also gives the sample
size: $192$ paired runs per point.
\item \emph{Local versus global Lipschitz constant (\S9.2).} At $\gamma=1.625$\,m/s
the choice of constant no longer changes which corridors certify $\theta=0.25$\,s
(the tube is $2.24$\,m with the global constant and $2.65$\,m with the local one).
The sentence claiming that it does was removed.
\item \emph{Wording errors.} ``Airspeeds below stall'' now reads ``above stall''.
The saturation sentence in \S7 had its direction reversed; it now says that
loosening the detector beyond $s$ does not lower the floor any further. Secure
state estimation assumes a bound on \emph{how many} sensors are attacked, not
knowledge of which ones. Verisig and NNV concern \emph{neural-network}-controlled
systems. Two cross-references to runtime assurance pointed to the Discussion,
which does not cover it, and now point to Related Work. The denial-duration
exponent is written $e^{LT_d}$. ``Lemma~1'' is now a live reference.
\end{itemize}

\textbf{Paper I.} The Interface-transfer statement is now an \emph{Assumption}
environment rather than a Definition. The sentence on ``two of the four factors''
in the certificate now names $L$ and $m$ explicitly. The slot-budget bound in the
introduction now matches the proposition ($k_s-1$ verification times). The
freshness table caption referred to a ``holdover column'', which does not exist;
it now refers to the note column.

% ---------------------------------------------------------------------
\section{The undetected-delay row of the pairing table}

The pairing table reported the measured undetected delay at $\theta=0.25$\,s as
$0.00$\,s. The cause was the campaign attacker, which draws its per-hop delays from
$U[1,7]$\,s or $U[1,10]$\,s whatever the threshold, so every malicious delay at
$\theta=0.25$\,s was flagged. We have now added a threshold-aware attacker
(\texttt{experiments/tuned\_attacker.py}) and run it on the same grid: three
topologies, 13 thresholds, 8 seeds and 4 routing policies, 1{,}248 runs in all. It
adds the largest delay that keeps each hop below $\theta$ after allowing for the
benign forwarding jitter of $0.05$--$0.18$\,s.
\begin{itemize}[leftmargin=1.4em,itemsep=2pt]
\item It is never flagged, at any threshold, so the theorem's premise of an
attacker that stays just below the threshold is now demonstrated rather than
assumed.
\item Its undetected delay never exceeds $H\min(\theta,s)$ in any run, so
Lemma~1 holds against an attacker that is trying to reach it.
\item At $\theta=0.25$\,s it realises $0.36$\,s (median; maximum $0.41$\,s) of the
$0.50$\,s budget, that is 73\%; the shortfall is the jitter it must leave room for.
Once the budget saturates it realises 99\%, with no undetected hop above $4.99$\,s,
so the cap at the slack is tight.
\item An attacker that injects $0.99\,\theta$ without the jitter allowance is
flagged on every hop below the slack, and above the slack every delay misses its
window.
\end{itemize}
The pairing table now has both rows (campaign attacker and threshold-aware
attacker), and \S6.1 and the worked example in \S9.3 report the result in two
sentences. The certificate still uses the worst-case budget, so no certified
number changes. With the measured $0.41$\,s, the tube at $\theta=0.25$\,s would be
$2.19$\,m, so the $2$\,m corridor fails on measurement as well as on the bound.

\end{document}
